\chapter{Evidence, Obligations, and Program}\label{evid:chapter}

\section{Reading a bounded receipt}\label{evid:reading}

Every chapter of this monograph has, at some point, deferred to ``the bounded Gate-4 receipt.'' This chapter pays that debt. It states the receipt exactly, explains what each of its numbers is and is not evidence for, walks through what the test suite actually exercises, and then does the complementary work that a candidate's evidence chapter owes its reviewers: a substantive treatment of each obligation that remains \emph{open} --- the six gaps between this kernel and any change of status. The chapter closes with related work and with the design-inversion thesis that the whole monograph has been arguing.

One framing point before the numbers. A receipt, in this program's vocabulary, is a historical record of software behavior: on a stated date, at a stated commit, a stated set of tests ran against stated artifacts and passed, and the record binds those statements together with digests. That is all it is. It is not an experiment (nothing empirical about cognition or CAS content was measured), not a review (no independent party has evaluated anything), and not authority (nothing in it licenses any claim or product use). The discipline of reading receipts this narrowly is not modesty for its own sake; it is the same discipline the kernel enforces mechanically in its constant-false authorization invariant, applied by the authors to their own evidence. A monograph about a kernel that cannot be talked into overclaiming should not itself overclaim, and this chapter is where that symmetry is most load-bearing.

\section{The bounded Gate-4 receipt, exactly}\label{evid:receipt}

\begin{finding}[Gate-4 receipt for CAS Core v8, 2026-08-04, commit \texttt{a1cd9901}]\label{evid:find:receipt}
At the bound receipt:
\begin{enumerate}
\item The focused CAS Core v8 suite records \emph{61 passing test methods}, with suite test-ID SHA-256
\begin{center}
\texttt{E2AF5014DE28DD0702211EF32A2757B99B869A13FAFBB0B8D1B18A50A9F8FEF5}.
\end{center}
\item The focused suite ran within an integrated three-suite run of \emph{208 test methods} and \emph{272 parameterized subtests} over \emph{54 hash-bound artifacts}.
\item The claim registry holds \emph{11 claim records}, all Ring~4 and \textsc{not licensed}.
\item The review matrix holds \emph{53 review criteria}, of which \emph{20 are open}.
\item The repository-wide collector was \emph{not clean}: four pre-existing package-path namespace collisions surfaced after 1{,}200 collected items. Consequently, \emph{no repository-wide verification claim is made}; the receipt covers the focused suites only.
\end{enumerate}
\end{finding}

Each line deserves unpacking, because each is a different \emph{kind} of statement.

\paragraph{The 61 methods and the suite test-ID hash.} The count of 61 is the least informative number on the receipt --- test counts are famously gameable and famously incomparable across projects --- and it is anchored by the number next to it, which is the informative one. The suite test-ID SHA-256 is a digest computed over the sorted, fully qualified identifiers of the tests in the focused suite. Its function is to make ``the suite passed'' a well-defined historical statement: the receipt binds not a count but an exact set, so a later run claiming continuity with this receipt must present the same test-ID digest, and any addition, removal, or renaming of tests changes the digest and visibly breaks the continuity claim. This is the kernel's own medicine applied to its evidence: just as a compatibility receipt binds compiler bytes rather than a compiler name, the suite receipt binds test identities rather than a test count. What the digest does \emph{not} do is certify test quality --- 61 trivial assertions would hash as stably as 61 sharp ones. Quality is addressed, to the extent it can be without independent review, by Section~\ref{evid:suite}'s account of what the fixtures actually exercise, and finally only by open obligation six.

\paragraph{The integrated run: 208 methods, 272 subtests, 54 hash-bound artifacts.} The focused suite did not run alone. The integrated run spans the Gate-4 lane's three suites --- Core and its companion-facing surfaces --- and the numbers say two distinct things. The 208/272 pair records method-level and parameterization-level granularity separately, because parameterized subtests (the same assertion body over a table of fixtures, heavily used in the encoding and loader rejection catalogs) inflate neither count deceptively: 208 methods exist as code; 272 test outcomes were recorded. The 54 hash-bound artifacts are the input side: source files, registry ledgers, schema files, compiler bytes, and fixture data whose SHA-256 digests are recorded in the run's manifest, so that the receipt's claim is closed under its inputs --- \emph{these} tests, over \emph{these} bytes, passed. An artifact drifting after the receipt does not falsify the receipt; it merely, and detectably, ends the receipt's applicability to the present. That is the intended semantics of all bounded evidence in this program: receipts age out; they are never retroactively wrong about their own run.

\paragraph{Eleven claim records, all Ring 4, all \textsc{not licensed}.} This line is about the registry's \emph{content} rather than the software's behavior, and its evidentiary value is its uniformity. The claim registry at the receipt contains no licensed claims, no ring other than 4, and no record whose boundary triple deviates from (\textsc{open}, \textsc{not licensed}, disclosure false). The kernel's loader enforces that malformed records cannot load; it cannot enforce that the program's humans refrain from \emph{well-formed} overclaims --- a syntactically valid record asserting a licensed status would load if its bytes were written. That no such record exists is therefore a fact about program discipline, checkable by anyone who reads the eleven records, and the receipt pins the ledger bytes under which the check holds.

\paragraph{Reading the eleven records.} The claim records' content, not just their uniform status, rewards inspection, because the registry demonstrates what Ring-4 claim-keeping looks like in practice. The records fall into three classes. \emph{Mechanism claims} assert that a named kernel mechanism implements its written contract --- for example, that the canonical encoder realizes the closed rule set --- each bound to the specification digest and the exercising fixture families; these are the claims the Gate-4 receipt is \emph{about}, and even they are registered \textsc{not licensed}, because tested-by-authors is not reviewed. \emph{Boundary claims} assert what the kernel does \emph{not} do --- no authorization path, no v8 semantics, no sandbox --- registered as claims precisely so that their negations cannot later be asserted as having been implicit; a program that registers its refusals makes its future overclaims checkable against its past. \emph{Bridge claims} cover the narrow guarantees of the integration surfaces (type-correct composition, marked covers proving nothing), each phrased with its scope limitation inside the claim text itself rather than in surrounding commentary. Nothing in the registry asserts empirical validity, cognitive-scientific content, or product fitness --- there is no record whose licensing would even be meaningful under the current evidence, which is itself a discipline exhibit: the registry contains only claims of the kind the program could someday support.

\paragraph{Fifty-three review criteria, twenty open.} The review matrix is the program's enumeration of what an independent reviewer would need to examine: per-mechanism specification-to-implementation correspondences, fixture adequacy questions, boundary-language audits, and the obligations of Section~\ref{evid:obligations} decomposed into checkable items. Thirty-three criteria are recorded as addressed \emph{by the authors' own lights}; twenty are open, and the open set includes every criterion that requires independence (no self-review can close it), every criterion blocked on unbuilt infrastructure (hermetic bundle, conformance campaign), and several fixture-coverage questions the authors judged honestly unresolved. Representative members of the open set, to give the ratio texture: whether the injective-rendering argument should be reconstructed machine-checked rather than accepted as a design argument with regression fixtures; whether the encoder's depth and node bounds are set at defensible values or merely at values (a policy criterion, not a code one); whether the lock discipline's constants deserve a contention study on the program's actual hosts; whether the fixture base for cross-namespace links covers all seven-choose-two namespace pairs or only the historically incident-bearing ones; and the standing criterion that every ``by construction'' phrase in the specification be traced to the constructor code that makes it true. The 20/53 ratio is not a progress bar and should not be read as ``62 percent done''; the open criteria are not smaller than the closed ones, and the closed ones reopen under independent review by design.

\paragraph{The collector that was not clean.} The last line of the receipt is its most important discipline exhibit. Running the repository-wide test collector --- the step that would justify a claim of the form ``the whole repository's suites pass'' --- surfaced four pre-existing package-path namespace collisions after 1,200 items were collected: directories in the wider monorepo whose package paths alias under the collector's import semantics, a defect of the surrounding repository, not of the Core suites. The honest options were to fix the wider repository (out of scope for a bounded candidate receipt), to exclude the colliding paths and claim the remainder (a silently narrowed claim wearing a broad one's language), or to record the collector as not clean and bound the claim to the focused suites. The receipt does the third. The irony is noted and intended: the failure that bounds the receipt is a \emph{namespace collision} --- the exact failure family this kernel exists to make loud. It was loud. The bound is the evidence that the program's evidence discipline operates on itself.

\subsection{Re-verification: what a reader can check today}\label{evid:reverify}

A bounded receipt is only worth its bound if a reader can, in principle, re-perform the binding. The re-verification procedure for Finding~\ref{evid:find:receipt} is mechanical and worth stating as a protocol, both because reviewers will need it and because writing it down exposes exactly which steps depend on trust that has not yet been earned.

\begin{protocol}[Receipt re-verification]\label{evid:prot:reverify}
Given the repository at commit \texttt{a1cd9901}:
\begin{enumerate}
\item Enumerate the focused suite's test identifiers by collection (without execution), sort them under the declared ordering, canonically encode the sorted list, and compute its SHA-256. Compare against \texttt{E2AF5014\ldots}: agreement establishes that the test \emph{set} is the receipted one.
\item Hash the 54 manifest artifacts and compare against the manifest's recorded digests: agreement establishes that the input bytes are the receipted ones.
\item Execute the focused suite and the integrated run; compare method and subtest counts and pass status.
\item Load the claim registry through the Core loader; verify it loads cleanly, contains 11 records, and that every record's ring and boundary fields read Ring~4, \textsc{open}, \textsc{not licensed}.
\item Run the repository-wide collector; observe that it reproduces the four package-path collisions after 1{,}200 items, confirming that the receipt's bound was forced, not chosen for convenience.
\end{enumerate}
\end{protocol}

Two honesty notes on this protocol. Step (3) re-executes on the reader's stack, not the receipt's: agreement is expected but is evidence about the reader's environment as much as the original, and disagreement would be a finding about environmental sensitivity, not necessarily about the receipt (this is obligation 1's shadow falling across even simple re-verification). And the whole protocol verifies \emph{internal} consistency --- receipt against repository --- not the repository against the world: a reader receives commit \texttt{a1cd9901} through channels (version control, filesystem) whose integrity the kernel records but cannot vouch for. Re-verification under these caveats is still far from empty: it reduces the trust surface to ``the bytes I was handed,'' which is exactly the reduction every mechanism in this kernel performs, applied once more at the outermost layer.

\begin{boundary}[What the receipt establishes]\label{evid:bd:receipt}
Finding~\ref{evid:find:receipt} establishes: that the identified test set, over the identified artifacts, passed on the identified date at the identified commit. It does not establish hermetic execution (open obligation 1), cross-implementation conformance (2), any v8 semantic capability (3), any claim authority (4), multi-file program coverage (5), or independently reviewed correctness (6). It establishes nothing repository-wide. All content remains Ring~4, \textsc{open}, \textsc{not licensed}.
\end{boundary}

\section{What the suite exercises}\label{evid:suite}

A test-ID digest tells a reviewer \emph{which} tests ran; this section tells them what the tests are \emph{about}, mechanism by mechanism, with emphasis on the suite's governing bias: negative fixtures dominate. A fail-closed kernel's specification is mostly a list of refusals, so most of its tests are demonstrations that refusal happens --- the suite spends more methods proving that wrong things are rejected than that right things are accepted, and the threat model of the previous chapter maps onto the fixture classes nearly one-to-one.

\paragraph{Reference identity and injective rendering.} Fixtures construct references across all seven namespaces and assert full-tuple structural equality semantics: equal tuples equal, any single-field perturbation (including \texttt{definition\_sha256} alone) unequal. Cross-namespace fixtures with identical local identifiers assert inequality (threat T1's mechanical core). The rendering fixtures are the delimiter gauntlet: components containing the join delimiter, empty components, components that are prefixes of one another, and pairs of distinct tuples crafted toward collision, each asserted to render distinctly and to round-trip through the parse (T2). Link fixtures assert that \textsc{translation}, \textsc{alignment}, and \textsc{hypothesis} construct and that no operation yields identity from any of them.

\paragraph{Boundary consistency and recheck.} Constructive fixtures assert that the Ring-4 triple builds and that every contradictory combination --- licensed product under \textsc{open} admission, disclosure true in Ring 4 --- fails at construction. Recheck fixtures exercise the derive-and-recheck path: operations receiving stale or mutated boundary context are shown to recompute rather than trust, and self-digest fixtures verify that mutation after construction is detected while asserting (as a tested property, not just prose) that a valid self-digest grants nothing (T7's first layer).

\paragraph{Canonical encoding.} The largest parameterized block. Acceptance fixtures pin exact expected bytes --- not just stable bytes --- for representative structures, so encoder drift against the specification is caught even if it drifts deterministically. Rejection fixtures cover the full catalog: sets and unordered containers, binary floats in every position, NaN, both infinities, untagged scalar ambiguity, cyclic structures, depth-limit and node-limit overflows (T4). The type-pun fixtures encode $\mathrm{Decimal}(\texttt{"1.0"})$, the string \texttt{"1.0"}, and the integer $1$ and assert pairwise distinct byte outputs, with the same pattern over further pun families. Newline fixtures encode fixture files with lone-CR, CRLF, and lone-LF content and assert that raw and comparison digests, and the three counts, are recorded unconflated (T5).

\paragraph{Compatibility protocol.} Hash-verification fixtures run the worker against correct and corrupted source and compiler bytes, asserting refusal on any digest mismatch (T6). Sanitization fixtures inspect the worker's recorded runtime for the \texttt{-I -S -B} profile and sanitized environment, and assert the runtime-binding fields are populated before-execution as specified. Lock fixtures exercise acquisition, bounded timeout behavior, and the no-retry discipline (T10); failure-taxonomy fixtures feed the worker malformed and empty transport responses and assert fail-closed classification with only exit code, byte counts, and digests recorded; divergence fixtures present mismatched dependency profiles and assert the bounded, hash-valued error shape --- both profile hashes, at most 32 keyed difference paths, no cleartext environment values (T11). Replay fixtures assert that the structural-decision path is unreachable without a fresh execution (T9).

\paragraph{Registry loading.} The second large parameterized block: one fixture family per rejection --- duplicate JSON keys, unknown root fields, unknown record fields, wrong enum values, wrong shapes, invalid identifiers, ambiguous versions, missing required metadata --- each asserting load failure with a diagnostic naming the defect (T8). Substitution fixtures attempt the constructor and dataclass-copy routes around the loader and assert rejection; immutability fixtures attempt post-load mutation; digest fixtures assert that selected-record digests change when ledger bytes change, binding records to their registry context.

\paragraph{Authorization invariants.} The smallest and most consequential block: fixtures assert that \texttt{authorization\_granted} and the compatibility spelling \texttt{allowed} are \emph{false} across every constructible decision outcome, including outcomes whose every structural field verifies and whose registry input says \textsc{licensed} --- the tested form of the statement that there is no code path to yes (T7's terminal layer). A static-search check asserts the absence of any assignment of true to either field in the kernel source.

\paragraph{Fixture philosophy.} Three commitments govern the fixture base, and stating them lets a reviewer audit the suite against its own standards rather than against generic testing folklore. First, \emph{negative dominance}: for every acceptance family there is at least one rejection family, and the rejection catalogs (encoding, loader) are exercised exhaustively --- every enumerated rejection reason has at least one fixture, because an untested refusal is a refusal that can silently stop happening. Second, \emph{byte-pinned oracles}: wherever the specification determines exact output (canonical encodings, rendered identifiers, digests), fixtures compare against pinned byte strings checked into the fixture set, not against the implementation's own round-trip --- a suite that asks the encoder whether the encoder is right tests only self-consistency, the failure mode that let deterministic-but-wrong drift pass in earlier program tooling. Third, \emph{adversarial provenance}: a substantial share of the negative fixtures are transcriptions of the threat model's scenarios (delimiter crafting, subclass impersonation, pun families, constructor bypass), so the suite is, in part, the threat model made executable. Table~\ref{evid:tab:blocks} records the correspondence.

\begin{table}[htbp]
\centering
\caption{Suite blocks against mechanisms and threats.}
\label{evid:tab:blocks}
\begin{tabular}{@{}lll@{}}
\toprule
Suite block & Mechanism chapter & Threats exercised \\
\midrule
Identity and rendering & semantic identity & T1, T2, T3 \\
Boundary & research boundary & T7 (recheck layer) \\
Canonical encoding & encoding & T4, T5 \\
Compatibility protocol & worker protocol, concurrency & T6, T9, T10, T11 \\
Registry loading & registry integrity & T8 \\
Authorization invariants & decision model & T7 (terminal layer) \\
\bottomrule
\end{tabular}
\end{table}

\paragraph{Coverage, honestly measured.} Line- and branch-coverage percentages are deliberately not headlined on the receipt, and the reason belongs in the record. For a fail-closed kernel, the coverage metric that matters is \emph{rejection coverage}: the fraction of enumerated refusal behaviors --- catalog entries, protocol failure classes, invalid boundary combinations --- that have at least one fixture demonstrating the refusal, with its diagnostic shape asserted. This metric is countable from the specification (the catalogs are finite and written down), it is the metric the parameterized blocks were built against, and at the receipt it is complete for the encoding and loader catalogs by construction --- each catalog entry is a fixture row. Conventional statement coverage is recorded in the run artifacts for reviewers but is not cited as evidence here, because in this codebase it is dominated by exactly the code whose correctness the metric cannot speak to: error-message formatting, record field plumbing, and other code that executes identically whether or not its semantics are right. A high statement-coverage number would decorate the receipt without strengthening it; the program's standing rule --- aggregate metrics never get to be the audit --- applies to its own test evidence first.

\begin{remark}[What the suite does not exercise]\label{evid:rem:notexercised}
For symmetry: the suite contains no performance tests (the kernel makes no latency claims), no concurrency stress beyond the lock discipline's specified behavior (no fairness claims exist to test), no fuzzing campaign (open review criterion), and no cross-platform matrix beyond the platforms in the integrated run's recorded environments. Each absence is a review criterion, several among the twenty open.
\end{remark}

\section{The six open obligations}\label{evid:obligations}

The obligations below are the candidate's honest distance from trustworthiness. Each is stated, then discussed substantively: why it is open, what closing it requires, and what misreading its openness invites.

\subsection{Obligation 1: hermetic subprocess and toolchain bundle}\label{evid:ob:hermetic}

Python-level isolation --- the \texttt{-I -S -B} worker with sanitized environment --- removes the inherited-state failure class but leaves the executing stack itself unverified: the interpreter binary, every native extension and shared library it loads, and the operating-system surface beneath them are recorded by name and version, not content-addressed. The gap is stated in every receipt as \texttt{false} capability fields, and the threat model locates the consequence exactly: an attacker or an upgrade at the native layer defeats the compiler-identity binding at a layer the protocol only describes.

Closing the obligation means a hermetic bundle: an execution environment in which the interpreter, its standard library, all native dependencies, and ideally the OS userland are identified by digest and assembled from those digests at execution time, in the spirit of reproducible-builds discipline \cite{reprobuilds}. The engineering is well-trodden in other ecosystems (content-addressed store, pinned toolchains), but the obligation is not merely to adopt a tool: the receipt schema must grow fields for the bundle identity, the replay requirement must be extended to demand bundle-identity agreement, and the failure taxonomy must classify bundle mismatches.

Two design questions must be answered under review before implementation, and they are recorded here because their answers shape the receipt schema. \emph{Granularity}: a bundle digest over one opaque image is easy to record and nearly useless to diagnose --- a mismatch says ``something differs'' --- while per-component digests (interpreter, each native library, each wheel) make divergence reports as diagnosable as the dependency-profile mechanism already is, at the cost of a manifest schema and a canonical component ordering; the kernel's existing preference for bounded, keyed difference reporting argues for the per-component design. \emph{Windows fidelity}: the program's primary hosts are Windows machines, where content-addressing the loader-visible surface (system DLLs, side-by-side assemblies) is materially harder than on POSIX; a bundle design that is hermetic on Linux and approximate on Windows must say so per-platform in its capability fields rather than reporting the strongest platform's property globally. The honest sequencing note: until this closes, obligation 2's conformance campaign is weaker than it looks, since cross-implementation agreement observed on unverified stacks conflates implementation behavior with environment behavior. The misreading to guard against: ``isolation'' in this monograph never means sandbox, and no current receipt is evidence against a host-level adversary.

\subsection{Obligation 2: Python/Rust conformance}\label{evid:ob:conformance}

The program's architecture places a Rust engine beside the Python kernel, and the seed specification permits a separately hash-bound native engine profile. None is admitted in v0.1: engine profiles remain separate, and no receipt asserts that the Python and Rust implementations agree on anything. The obligation is a differential-testing campaign: a shared corpus of inputs (encoding structures, reference tuples, registry byte fixtures, compiler sources) executed against both implementations, with canonical outputs compared at the digest level, counterexamples minimized and receipted, and the corpus itself hash-bound and grown adversarially --- generative fuzzing biased toward the rejection catalogs, where implementations most plausibly diverge (one accepts what the other refuses).

The reason this is an obligation rather than a task is its acceptance criterion. ``The corpus passed'' is bounded evidence about the corpus; the conformance claim wanted is universally quantified, and the campaign can only ever approach it. The program's discipline for this is the same as everywhere else: conformance receipts will be bounded statements naming their corpus digest, and the phrase ``Python/Rust conformant'' without a corpus identity attached will remain unlicensed language in the lane.

The campaign's hardest technical content is not comparison but \emph{specification arbitration}. When the implementations disagree, three verdicts are possible --- Python wrong, Rust wrong, specification ambiguous --- and the third is the valuable one: every ambiguity found is a place where the written contract underdetermines behavior and where a future third implementation would diverge again. The campaign must therefore feed a specification-erratum stream, each erratum a reviewed tightening of the contract text with its motivating counterexample attached, and the contract-versioning rules of the integration chapter govern how tightened text propagates to companions. There is also a floating-point-shaped trap to name in advance: the canonical encoder's rejection of binary floats means the conformance corpus never exercises float formatting --- by design --- but Rust-side arithmetic feeding tagged decimals must be checked for exactness at the construction boundary, since a Rust path that computes in \texttt{f64} and converts at the edge would reintroduce, invisibly, the representability drift the encoder exists to exclude. Until the campaign exists, the operative rule is stark: no evidence produced through one implementation may be cited as evidence about the other.

\subsection{Obligation 3: no v8 parser or semantic migration}\label{evid:ob:parser}

The only admitted compatibility profile is \texttt{legacy\_only}: it executes the hash-verified legacy Python compiler for comparison purposes and produces no CAS v8 semantic translation. There is no v8 parser, no v8 semantics, and no migration path from v7.7 artifacts to any v8 form --- and this is a design decision, not a delay. The kernel was built first precisely so that when v8 semantics arrive, they arrive into an environment where their sources are hash-identified, their evidence encodings are canonical, their claims are registered Ring-4, and their execution is receipted.

The obligation's substance is therefore mostly \emph{criteria}: any future v8 parser must arrive with its grammar as a reviewed artifact, its conformance to the legacy compiler on the shared subset established through the compatibility profile (which exists exactly to make such comparisons receiptable), and its divergences from legacy behavior enumerated as reviewed decisions rather than discovered as surprises.

A word on why migration, specifically, is refused rather than merely deferred. A migration tool is a semantic claim generator: every artifact it converts carries an implicit assertion that the v8 form means what the v7.7 form meant, multiplied across the entire migrated corpus at once. Shipping such a tool before v8 semantics are reviewed would manufacture, in bulk, exactly the unlicensed equivalence claims that Rule-level discipline forbids one at a time --- the no-silent-identity rule applied at corpus scale. When migration does arrive, the kernel's shape for it already exists: each conversion should be receipted like a compatibility execution (source identity in, tool identity, target identity out, boundary attached), so that a migrated corpus is a set of auditable derivation records \cite{prov} rather than a fait accompli. The misreading to guard against is treating the kernel's existence as progress on v8 semantics. It is progress on being \emph{able to test} v8 semantics; the semantic work --- the hard, domain-laden work of saying what v8 notation means --- has not begun, and no artifact in this monograph moves it.

\subsection{Obligation 4: no claim authority}\label{evid:ob:authority}

The constant-false invariant is the kernel's signature, and it is also, read as program planning, a deferred liability: a research program eventually needs to make claims, and the current kernel can only refuse them. The future design's minimum constraints are already stated in the runtime decision model: an exact trusted registry and schema manifest; an authenticated authority and trust root; reviewer and admission receipts; Evidence Packet identities and verification status; claim ID and exact version; and the code, model, environment, semantic, and calibration identities of the requested operation.

The discussion this obligation deserves is about \emph{difficulty class}. Every other obligation is engineering; this one is institution-building. A trust root is only meaningful if key custody, reviewer authentication, admission process, and revocation are real practices with real people bound to them --- the cryptography is the easy tenth of the problem \cite{anderson}. The threat model is explicit that its current T7 analysis (residual risk zero, because ``yes'' is unrepresentable) expires the day this design lands, and that the full laundering threat family must be re-analyzed against the new surface.

The program's specific vulnerability here should be named rather than implied: it is, at present, substantially a single-operator program assisted by AI tooling, and the historical self-hashed licensing receipts were produced inside exactly that configuration. An authorization design whose ``authenticated reviewer'' could be instantiated by the same operator who requests the claim, or by tooling acting on the operator's behalf, would be ceremony, not authority --- constant-false with extra steps, minus the honesty. The minimum-constraints list therefore implies an organizational fact the cryptography cannot supply: the reviewer role must be filled by parties independent of the program before the role's credentials mean anything, which ties this obligation to obligation 6 not just in sequence but in substance. Revocation deserves equal notice: claims, once authorized, must be revocable when their underlying evidence ages out (a receipt superseded, a review withdrawn), so the design is not ``grant'' but a full lifecycle --- grant, cite, re-verify, revoke --- and the registry schema must represent all four states without ambiguity. The program's stated posture is that this obligation closes \emph{last} among the mechanism obligations: authorization built before hermetic execution and independent review would authorize on evidence the other obligations say is not yet trustworthy. Until then, the kernel's ``no'' is not a placeholder; it is the correct answer to every claim-bearing request made of an unreviewed system.

\subsection{Obligation 5: no import-DAG manifest}\label{evid:ob:importdag}

Protocol step (3) rejects \texttt{use} imports: a compatibility execution covers exactly one source file, and multi-file CAS programs are outside the admitted profile. The missing piece is an import-DAG manifest: a canonical, hash-bound description of a program's full dependency graph --- every imported source identified by digest, the graph's shape recorded, cycles rejected or declared, and the whole manifest itself canonically encoded so that ``the program'' as an evidence object means the closed set of bytes, not a root file plus whatever the filesystem happened to resolve.

The reason to refuse partial solutions is the threat model's T6: resolving imports at execution time from an unmanifested filesystem reintroduces, at the program level, exactly the identity gap the protocol closed at the compiler level --- the receipt would bind the root file's bytes while the semantics depended on unbound neighbors. Rejecting \texttt{use} outright, rather than resolving-and-recording opportunistically, was the deliberate v0.1 choice: an opportunistic recorder would have produced receipts whose completeness varied with filesystem state, and a receipt schema with optional completeness is a receipt schema whose consumers must each decide, privately, what a receipt means --- the interpretive scatter the kernel exists to end. The manifest design also has one genuinely open semantic question: whether the DAG's identity should bind resolution \emph{order} (making the manifest sensitive to traversal, hence stricter) or only the digest \emph{set and edge relation} (making semantically irrelevant orderings receipt-equivalent). The kernel's general bias toward stricter identity says bind the order; the conformance obligation's interest in cross-implementation stability says bind the relation; the tension is recorded for review rather than resolved here. The obligation interacts with obligation 1 (a manifest of source digests still executes on an unverified stack) and with obligation 3 (a v8 parser will need module semantics, and the manifest schema should be designed once, for both). Until it closes, the honest statement is that compatibility evidence exists for single-file programs only, and the suite's fixtures respect that bound.

\subsection{Obligation 6: independent review}\label{evid:ob:review}

The last obligation subsumes the rest: independent software review, formal-methods review of the specifications and their invariant arguments, semantics review of the CAS content the kernel will eventually carry, and product-safety review of any use beyond research --- all before any status change. Twenty open criteria in the review matrix are the concrete docket, and several closed ones are closed only self-attestedly and reopen under independence.

Two design features exist specifically to make this obligation dischargeable rather than aspirational. First, everything reviewable is an artifact: specifications, source, fixtures, receipts, ledgers, and this monograph are byte-identified, so a reviewer can bind their findings to exact digests and their review becomes, itself, a receipt in the program's sense --- which is the shape the future trust root (obligation 4) will need review evidence to take. Second, the strict dependency direction of the integration chapter caps the review surface: Core is reviewable in isolation, and its review does not transitively obligate a reviewer to the companions' domain content.

The four review disciplines are not interchangeable, and conflating them is itself a review failure to preempt. \emph{Software review} asks whether the implementation matches the contracts and the fixtures exercise them --- the twenty open criteria's largest share. \emph{Formal-methods review} asks whether the invariant arguments (injective rendering, encoding determinism, loader completeness against its rejection catalog) actually hold as stated, ideally by independent reconstruction or machine-checked formulation; the monograph's design arguments are proofs of intent, not verified proofs. \emph{Semantics review} concerns the CAS and TKS content the kernel will eventually carry and cannot begin in earnest until obligation 3 produces reviewable semantic artifacts; its present scope is confined to checking that the kernel makes no semantic commitments prematurely --- a review of absences. \emph{Product-safety review} is contingent: it activates only if any use beyond research is ever proposed, and the boundary fields exist precisely so that its trigger condition is observable. A reviewer engaged for one discipline should be told which one, in these terms. What the program cannot manufacture is the independence itself; no internal process substitutes, and the boundary language of every artifact --- \textsc{draft / open / ring 4 / not licensed / held} --- remains in force until it arrives.

\section{Program: sequencing the obligations}\label{evid:program}

The obligations are not a list but a partial order, and stating the order is part of the evidence program. Obligation 1 (hermetic bundle) is upstream of the strong form of 2 (conformance observed on verified stacks) and strengthens every receipt retroactively issued after it. Obligation 5 (import-DAG manifest) is upstream of any serious use of 3 (a v8 parser tested only on single files would be tested on toys). Obligation 4 (authority) is explicitly sequenced last among mechanisms, for the reason given in its discussion. Obligation 6 (independent review) is not a stage but a gate crossing all of them: each closed mechanism obligation produces artifacts that enter review, and no closure is final until reviewed. The program's near-term work is therefore 1 and 5 in parallel --- both are bounded engineering with receipt-schema consequences that should be designed together --- with 2's corpus construction beginning against the current stacks (weak-form conformance now, upgraded to strong form when 1 lands), and 3 and 4 deliberately idle until their prerequisites exist.

This ordering is itself Ring-4 planning content: it binds no one, and independent review (which may reorder it) outranks it. It is recorded so that the program's future choices can be audited against its present reasoning --- the same function every other artifact in the lane serves.

\paragraph{Program risks.} A program section that lists only tasks is planning theater; the risks to the program deserve equal standing. Three are material. \emph{Erosion under friction}: the threat model's compound analysis already identified the pattern --- fail-closed refusals generate operational friction, and friction generates pressure to relax gates; the mitigation is structural (gates are protocol, not policy) but the pressure will fall on the review process that approves Core revisions, which is exactly where obligation 6's independence matters. \emph{Single-maintainer fragility}: the same organizational fact that complicates obligation 4 threatens continuity generally; the mitigation on record is that every artifact is byte-identified and self-describing enough for a successor or reviewer to reconstruct state from the repository alone --- a property this monograph is itself intended to serve. \emph{Scope seduction}: the kernel's discipline makes it tempting to route ever more program content through it, and each new consumer arrives with pressure to widen contracts (new encoder types, new namespaces); the versioning rules make widening a review event, but the cumulative drift of many small reviewed widenings is a risk no single review sees, and the periodic whole-surface re-review implied by obligation 6 is the only standing answer. None of these risks is retired by any receipt; they are properties of the program, not the software, and are recorded here so their management can be audited like everything else.

\section{Related work}\label{evid:related}

The kernel's mechanisms adapt established practice rather than claiming novelty, and the honest positioning is per-mechanism. The canonical evidence encoding is in the tradition of JSON canonicalization for stable digests, RFC 8785 being the nearest reference point \cite{jcs}; the kernel's departures --- rejecting floats outright rather than canonicalizing their formatting, tagging exact decimals, refusing unordered containers instead of sorting them --- trade generality for a smaller argument surface, which is the correct trade for an evidence encoder and would be the wrong one for an interchange format. The identity and receipt discipline --- content-addressing sources, toolchains, and artifacts; treating a build or run as trustworthy only relative to pinned input digests --- is the reproducible-builds program \cite{reprobuilds} applied to a research kernel, with the hermetic gap between the two stated as obligation 1 rather than elided. The separation of entity, activity, and derivation in receipts and traces follows the provenance-modeling tradition standardized in PROV \cite{prov}, restricted here to flat, canonically encoded records because the kernel refuses graph semantics it cannot yet bound. The threat-model method is Shostack's \cite{shostack}; the governing security-engineering stance --- integrity mechanisms are not authorization, trust roots are institutional before they are cryptographic, systems fail at seams --- is Anderson's \cite{anderson}, and the constant-false authorization invariant can be read as that stance taken to its fixed point: when no adequate trust root exists, the correct implementation of authorization is the constant function \emph{no}.

Three contrasts sharpen the positioning. Against \emph{interchange canonicalization}: RFC 8785 exists so that arbitrary JSON documents produced anywhere can be canonically compared, and it therefore must handle the full JSON value space including IEEE floats \cite{jcs}; the kernel's encoder serves the narrower purpose of digesting evidence it also \emph{produces}, and can therefore shrink the value space to what admits an exactness argument --- the two designs are correct for their respective problems, and importing either into the other's setting would be a defect. Against \emph{supply-chain attestation practice}: the reproducible-builds movement's central lesson is that verifiability must be designed in rather than audited on \cite{reprobuilds}, and its central sociological observation --- that reproducibility claims decay unless continuously re-exercised --- is the direct ancestor of the kernel's replay requirement; where the kernel departs is in refusing the ecosystem's common comfort of ``attestation documents'' as endpoints, since the program's history shows attestation-shaped artifacts being read as authority (threat T7). Against \emph{provenance modeling}: PROV provides a graph vocabulary rich enough to express agency, delegation, and derivation chains \cite{prov}; the kernel deliberately uses only its flattest fragment --- which entity, which activity, which inputs --- because agency and delegation are exactly the notions that must wait for an authenticated trust root, and a vocabulary that can express ``authorized by'' before the system can verify it is a laundering surface. In each contrast the kernel's choice is the narrower one, and the narrowness is the point.

Against this background, the kernel's contribution is compositional and contextual: a fail-closed assembly of these known disciplines into one substrate, for a research program with \emph{documented} drift and overclaim incidents, with the claim discipline enforced in the object language (unrepresentable states, constant-false fields) rather than in surrounding process documents alone. Whether that composition is itself sound is precisely what open obligation 6 exists to determine.

\section{Conclusion: the design inversion}\label{evid:conclusion}

The thesis of this monograph can now be restated with all of its supporting structure in view.

Conventional practice builds capable systems and then subtracts danger: add semantics, then validate; add authority, then audit; add integration, then review. The program's own history is a record of what that ordering costs when the system's outputs are claims --- positional correspondences hardened into identities before anyone checked, and structurally perfect receipts read as licenses because nothing in their design refused the reading. CAS Core 8 inverts the order. It begins with a kernel in which the historically observed failures are unrepresentable or fail closed --- identity is structural and namespaced, correspondences carry their kinds, encodings are closed and deterministic, byte identities are dual and unconflated, execution is hash-bound and replayed, registries are byte-backed, and authorization is the constant \emph{no} --- and it requires everything a future CAS v8 wants to assert to be \emph{added}, through gates that are themselves reviewed: a catalog for every digit, a manifest for every import, a trust root for every claim, a reviewer for every closure.

The evidence for the inversion's viability is deliberately modest and exactly bounded: 61 focused methods under one test-ID digest, inside a 208-method, 272-subtest, 54-artifact run; eleven claims, all Ring 4; twenty criteria open; one collector honestly not clean. The evidence for the inversion's \emph{necessity} is the program's incident history, and the evidence for its \emph{character} is distributed through every mechanism chapter: at each design fork, the kernel chose the option that made a failure loud over the option that made a workflow smooth.

The inversion has a cost model, and honesty requires stating it. Fail-closed kernels are slower to build against, hostile to prototyping, and indifferent to deadlines; every companion candidate paid integration friction that a permissive substrate would have waived, and the six obligations exist \emph{because} the kernel refuses to paper over them --- a permissive design would simply not have surfaced most of them as obligations at all. The wager of the whole program is that this cost is the cheap side of the ledger: that for software whose product is claims, the expensive failure is not the blocked workflow but the plausible artifact --- the conflation that survives review because nothing typed it, the receipt that becomes a license because nothing refused the reading. The program has paid the expensive failure before. The kernel is what it bought with the lesson: subtraction replaced by gated addition, and a substrate that will keep saying \emph{no} until closing its obligations earns it the right to say anything else.

\begin{boundary}[Final status statement]\label{evid:bd:final}
This monograph documents a Ring-4 research candidate: admission \textsc{open}, product \textsc{not licensed}, public disclosure withheld. Its receipts are historical software evidence under exact fixtures and artifacts. Nothing herein is an experiment, a certification, a review, a product representation, or a licensed claim, and no reading of this document's completeness, internal consistency, or test evidence converts it into one.
\end{boundary}
